%% kapa: exam=1 round=28 year=2017 period=1 form=A subject=civil_law subject_ko=민법 seq=27 number=27
%% source: Q-Net 감정평가사 기출문제 (cst00309) artlSeq=5199603 file=extracted/1008/5199603/2017년 제28회 감정평가사 1차시험 문제지/1교시 A형.pdf
%% conversion: kordoc 4.24.51 -> tools/gen1.py (auto); review=proofread-sonnet; flags=none
\begin{question}{27}
점유자와 회복자의 관계에 관한 설명으로 옳지 않은 것은? (다툼이 있으면 판례에 따름)
\begin{choices}
  \item 선의의 점유자가 과실을 취득한 범위에서는 그 이득을 반환할 의무가 없다.
  \item 유효한 도급계약에 기하여 수급인이 도급인으로부터 제3자 소유 물건을 이전받아 수리를 마친 경우, 원칙적으로 수급인은 소유자에 대하여 비용상환청구권을 행사할 수 있다.
  \item 악의의 점유자도 원칙적으로 필요비 전부의 상환을 청구할 수 있다.
  \item 점유물이 점유자의 책임 있는 사유로 멸실 또는 훼손된 경우, 악의의 점유자는 자주점유라도 손해 전부를 배상할 책임이 있다.
  \item 점유자가 과실을 취득한 경우에는 통상의 필요비는 청구하지 못한다.
\end{choices}
\end{question}
\answer{2}
